\chapter{Статус динамического замыкания спектрального рождения}
\label{ch:v6-spectral-transition-dynamic-closure-status-gate}

\section{Задача итогового гейта}

Ветвь получила строгий язык спектрального кванта перехода, но каждый
последующий тест отделял классификацию от физического рождения. Настоящий
гейт сводит результаты в один ledger и запрещает продолжать неизменённую
архитектуру там, где отсутствует общий динамический родитель.

\section{Закрытые уровни}

На уровне классификации закрыты:
\begin{enumerate}
 \item конечный носитель $\mathcal H_{15}$ и его Real/KO6-типизация;
 \item тёплицев класс $15=12+3$ и компонентные веса;
 \item локальный ковариант $W_\nu=\widetilde H\widetilde H^\dagger$ со
 сменой ранга $0\to1$ при $H=0$;
 \item физический сфалеронный поток $(3,1)$ и правила
 $\Delta B=\Delta L$, $\Delta(B-L)=0$.
\end{enumerate}
Это классификационное и кинематическое замыкание, а не теория рождения.

\section{Частично закрытый переходный уровень}

Полный электрослабый gauge--Higgs сектор допускает конечное седло с
$H(0)=0$: сфалерон \cite{dynamic-status-klinkhamer-manton1984}.
Он подтверждает физическую возможность перехода через нулевую моду, но
не реализует проектный класс $15$, имеет отрицательную моду и не является
устойчивым конечным объектом. Поэтому уровень существования контрольного
седла закрыт только как \emph{control}, а проектное седло остаётся открытым.

\section{Открытые динамические уровни}

В текущей версии отсутствуют:
\begin{enumerate}
 \item самозапускающаяся неустойчивость одного общего проектного действия;
 \item bounce или иной конечный квантовый седловой путь того же действия;
 \item мера траекторий и prefactor, необходимые для скорости
 \cite{dynamic-status-coleman1977,dynamic-status-callan-coleman1977};
 \item устойчивый локализованный фермионный endpoint;
 \item независимое численное предсказание массы, времени жизни или скорости.
\end{enumerate}

\begin{theorem}[Статус версии VI]
Версия VI закрывает классификацию, операторную кинематику и компонентные
selection rules спектрального перехода. Она не закрывает динамическое
рождение устойчивой наблюдаемой материи: не выведены проектное седло,
эндогенный trigger, вероятность перехода, скорость и устойчивый endpoint.
\end{theorem}

\section{Сводный ledger}

\begin{center}
\begin{tabular}{p{0.39\textwidth}p{0.19\textwidth}p{0.31\textwidth}}
\toprule
Уровень & Статус & Сертификат \\
\midrule
Классификация KO6/Toeplitz & закрыт & $15=12+3$ \\
Кинематика смены ранга & закрыт & $W_\nu:0\to1$ \\
Переходное седло & контрольный & стандартный сфалерон \\
Новый численный observable & частичный & только selection rules \\
Эндогенный запуск & открыт & общего действия нет \\
Вероятность и скорость & открыт & меры и prefactor нет \\
Устойчивый endpoint & открыт & решения и гессиана нет \\
\bottomrule
\end{tabular}
\end{center}

\section{Стоп-критерий неизменённой архитектуры}

\begin{nogocriterion}
После независимого отсутствия общего седла, trigger, скорости и устойчивого
endpoint нельзя добавлять новые толкования рангов, индексов или весов и
называть это продвижением динамики. Продолжение допустимо только после
явного объявления новой модели с одним родительским функционалом.
\end{nogocriterion}

Минимальный контракт новой модели должен требовать:
\begin{enumerate}
 \item gauge- и Real-ковариантный путь на полном физическом носителе;
 \item одно действие для вакуума, перехода и конечного состояния;
 \item выведенную отрицательную моду или квантовое нуклеационное седло;
 \item вычислимую скорость с мерой и prefactor;
 \item устойчивый локализованный endpoint;
 \item хотя бы один новый численный observable, не внесённый как параметр.
\end{enumerate}

\section{Следующий гейт}

Следующий гейт
\path{version6_spectral_transition_new_model_minimal_requirements_gate}
должен превратить этот список в формальные приёмочные тесты. Нулевая
гипотеза: ни одна новая архитектура не допускается, пока не проходит все
обязательные тесты. Стоп-критерий: отсутствие хотя бы одного обязательного
пункта фиксирует не частичное рождение, а недопущенную модель.

\begin{protocol}
\begin{enumerate}
 \item закрытых уровней: \textbf{2};
 \item частичных или контрольных уровней: \textbf{2};
 \item открытых уровней: \textbf{3};
 \item классификационное замыкание: \textbf{да};
 \item кинематическое замыкание: \textbf{да};
 \item проектное переходное седло: \textbf{нет};
 \item эндогенный trigger: \textbf{нет};
 \item скорость рождения: \textbf{нет};
 \item устойчивый endpoint: \textbf{нет};
 \item динамическое замыкание: \textbf{нет};
 \item продолжение без новой модели: \textbf{остановлено}.
\end{enumerate}
\end{protocol}

Машинный сертификат сохранён в
\path{s2t_v6_spectral_transition_dynamic_closure_status_gate_results.json}.

\begin{thebibliography}{9}
\bibitem{dynamic-status-klinkhamer-manton1984}
F.~R.~Klinkhamer, N.~S.~Manton,
\emph{A Saddle-Point Solution in the Weinberg--Salam Theory},
Physical Review D 30 (1984), 2212--2220.

\bibitem{dynamic-status-coleman1977}
S.~Coleman, \emph{The Fate of the False Vacuum. I. Semiclassical Theory},
Physical Review D 15 (1977), 2929--2936.

\bibitem{dynamic-status-callan-coleman1977}
C.~G.~Callan, S.~Coleman,
\emph{The Fate of the False Vacuum. II. First Quantum Corrections},
Physical Review D 16 (1977), 1762--1768.
\end{thebibliography}